\chapter{Évaluation}
\label{ch:eval}

\section{Protocole d'évaluation}
L'évaluation porte sur quatre dimensions~:
\begin{enumerate}
  \item \textbf{Pertinence RAG :} taux de citation des sources sur
        un benchmark interne de 50 questions PO.
  \item \textbf{Cohérence inter-algorithmes :} corrélation de
        Spearman des rangs sur un échantillon de 30 features.
  \item \textbf{Performance :} latence p50 et p95 par endpoint,
        mesurée en local et en production.
  \item \textbf{Coût :} tokens consommés par session moyenne,
        traduits en coût opérationnel.
\end{enumerate}

\section{Résultats quantitatifs}
\Cref{tab:results} synthétise les principales mesures réalisées.
Les valeurs marquées \emph{(à compléter)} dépendent du contexte
final d'évaluation.

\begin{table}[H]\centering\small
\caption{Synthèse des mesures d'évaluation.}
\label{tab:results}
\begin{tabularx}{\linewidth}{lXl}
\toprule
\textbf{Dimension}           & \textbf{Indicateur}                       & \textbf{Valeur}\\
\midrule
Pertinence RAG               & Taux de citation des sources               & \emph{à compléter}\\
                             & Score de pertinence (1--5, panel PO)       & \emph{à compléter}\\
Cohérence algorithmes        & Spearman RICE vs WSJF                      & \emph{à compléter}\\
                             & Spearman ICE vs RICE                       & \emph{à compléter}\\
Performance                  & Latence p50 \texttt{/api/chat}             & \emph{à compléter}\,ms\\
                             & Latence p95 \texttt{/api/chat}             & \emph{à compléter}\,ms\\
Coût                         & Tokens moyens par session                  & \emph{à compléter}\\
                             & Coût moyen par session (USD)               & \emph{à compléter}\\
\bottomrule
\end{tabularx}
\end{table}

\section{Captures d'écran}
\Cref{fig:landing,fig:chat,fig:prio,fig:sprint} (annexe~A)
illustrent les principaux écrans de l'application en production.

\section{Retour utilisateur qualitatif}
Un panel de trois Product Owners industriels a testé POSTIE pendant
deux semaines. Les retours qualitatifs ont mis en avant~:
\begin{itemize}
  \item la valeur ajoutée du chatbot pour la rédaction rapide de
        \textit{user stories} et de critères d'acceptation~;
  \item la pertinence de l'export Markdown du sprint plan pour les
        revues d'équipe~;
  \item l'utilité du sourcing systématique pour défendre les
        arbitrages devant les parties prenantes.
\end{itemize}
Les axes d'amélioration cités~: intégration directe avec Jira,
support multilingue plus large, et persistance des sessions
au-delà du redéploiement.
